\chapter{Méthode de estimation}

\section{Méthode des moments}

\smallskip

\begin{mdframed}
\begin{definition}[Estimateur par la méthode des moments]

\begin{itemize}
    \item Soit un échantillon $X_1, \ldots, X_n$ de loi $\mathbb{P}_\theta$, $\theta \in \Theta$.
    \item Soit $\phi$ une fonction telle que $\mathbb{E}_\theta\left[\|\phi(X_1)\|\right] < +\infty$.
\end{itemize}
On souhaite estimer
\[
g(\theta) = \mathbb{E}_\theta[\phi(X_1)].
\]
On appelle \textbf{estimateur par la méthode des moments} :
\[
\hat{g}_n = \frac{1}{n} \sum_{i=1}^n \phi(X_i).
\]
\end{definition}
\end{mdframed}


\begin{notation}
\begin{itemize}
    \item $\hat{g}_n$ est consistant (fortement) par la loi des grands nombres.
    \item $\hat{g}_n$ est sans biais (par linéarité de l'espérance).
    \item Plus généralement, si $\hat{g}_n$ est construit par la méthode des moments comme ci-dessus, et si $h$ est une fonction, on dira aussi que
    \[
    h(\hat{g}_n)
    \]
    est un estimateur par la \textbf{méthode des moments} pour l'estimation de $h(g(\theta))$.
\end{itemize}
\end{notation}


\begin{remarque}
À priori, on ne connaît pas les propriétés de cet estimateur (biais, variance, normalité asymptotique). Il faudra étudier ces propriétés au cas par cas, éventuellement en utilisant la méthode delta.
\end{remarque}

\begin{exemple}
\begin{enumerate}
    \item \textbf{Probabilité d'un événement}
    
    Soit $A$ un événement. On souhaite estimer
    \[
    g(\theta) = \mathbb{P}_\theta(X_1 \in A) = \mathbb{E}_\theta[\mathbb{1}_{X_1 \in A}],
    \]
    où
    \[
    \mathbb{1}_{X_1 \in A} = \begin{cases}
    1 & \text{si } X_1 \in A, \\
    0 & \text{sinon}.
    \end{cases}
    \]
    
    Par la méthode des moments :
    \[
    \hat{g}_{A,n} = \frac{1}{n} \sum_{i=1}^n \mathbb{1}_{X_i \in A}.
    \]
    
    Cet estimateur est la fréquence empirique de l'événement $A$.
    
    \item \textbf{Loi uniforme}
    
    Soit $X_i$ pour $i \in [ 1, n ]$ de loi $\mathcal{U}([\theta-1 , \theta+1])$ avec $\theta \in \mathbb{R}$.
    
    Alors
    \[
    \mathbb{E}_\theta[X_1] = \theta = g(\theta).
    \]
    
    Par la méthode des moments :
    \[
    \hat{g}_n = \overline{X}_n = \frac{1}{n} \sum_{i=1}^n X_i,
    \]
\end{enumerate}
\end{exemple}




\section{Méthode du maximum de vraisemblance}

On se place dans le cadre de l'estimation de $g(\theta)=\theta$.

\smallskip

\begin{mdframed}
\begin{definition}
On appelle vraisemblance associée à l'observation $X$ la fonction

$$
L : 
\begin{cases}
\Theta \to \mathbb{R} \\
\theta \mapsto L(\theta,X)
\end{cases}
$$

définie par

\begin{itemize}
\item Si $X$ v.a. discrète

$$
L(\theta,X) = P_\theta (X) = \mathbb{P}_\theta (X=x)
$$
\item Si $X$ v.a. à densité

$$
L(\theta,X) = f_\theta(X)
$$

où $f_\theta$ est la densité de $X$ sous $\mathbb{P}_\theta$
\end{itemize}
\end{definition}
\end{mdframed}

\begin{remarque}
Si $X=(X_1,\cdots,X_n)$ est un $n$-échantillon i.i.d.

\begin{itemize}
\item $X_i$ discret :

$$
L(\theta,X) = \prod_{i=1}^{n} P_\theta (X_i)
$$
\item $X_i$ à densité :

$$
L(\theta,X) = \prod_{i=1}^{n} f_\theta(X_i)
$$
\end{itemize}
\end{remarque}

\begin{exemple}
On considère $X \sim \mathcal{B}(15,p)$ avec $p$ inconnue. On observe $x=5$.

La fonction de vraisemblance est

$$
L(p,5) = \mathbb{P}_p (X=5) = \binom{15}{5} p^5 (1-p)^{10}
$$

C'est la probabilité d'observer $5$ quand le paramètre est $p$.

\begin{center}
\begin{tabular}{|c|c|c|c|c|c|}
\hline
$p$ & 0,1 & 0,2 & 0,3 & 0,4 & 0,5 \\
\hline
$L(p,5)$ & 0,01 & 0,1 & 0,21 & 0,19 & 0,09 \\
\hline
\end{tabular}
\end{center}

Quand $p=0,3$ il y a $21$ chances sur $100$ d'observer $x=5$.

Il est donc plus « vraisemblable » que le vrai $p$ soit égal à $0,3$ plutôt qu'à $0,1$.

En suivant ce raisonnement, la valeur de $p$ la plus vraisemblable est celle qui maximise la probabilité d'observer $5$.

C'est donc la valeur qui maximise la vraisemblance.

Pour trouver ce maximum, on utilise la log-vraisemblance :

\begin{align*}
l(p,x) &= \ln L(p,x) = \ln \binom{15}{x} + x\ln(p) + (15-x)\ln(1-p)
\end{align*}

On dérive par rapport à $p$ :

\begin{align*}
\frac{\partial l}{\partial p}(p,x)
&= \frac{x}{p} - \frac{15-x}{1-p}\\
&= \frac{x(1-p) - (15-x)p}{p(1-p)}\\
&= \frac{x-15p}{p(1-p)}
\end{align*}

Cette dérivée s'annule en $p=\frac{x}{15}$.

On vérifie qu'il s'agit bien d'un maximum :

$$
\frac{\partial^2 l}{\partial p^2}(p,x) = -\frac{x}{p^2} - \frac{15-x}{(1-p)^2} < 0
$$

Ici avec $x=5$ on a $\hat{p}=\frac{5}{15}=\frac{1}{3}$.

La valeur de $p$ la plus vraisemblable au vu de l'observation $x=5$ est $p=\frac{1}{3}$.
\end{exemple}

\smallskip

\begin{mdframed}
\begin{definition}
On appelle estimateur du maximum de vraisemblance (E.M.V.) de $\theta$ toute quantité $\hat{\theta}$ satisfaisant :

$$
\hat{\theta}^{\text{EMV}} \in \arg \max_{\theta \in \Theta} L(\theta,X)
$$

Souvent on considère la log-vraisemblance :

$$
\hat{\theta} \in \arg \max_{\theta \in \Theta} l(\theta,X)
$$

où $l(\theta,X) = \ln L(\theta,X)$.
\end{definition}
\end{mdframed}

\begin{proposition}[Méthode pratique de calcul]
Si $L(\theta, X)$ (ou $l(\theta,X)$) est dérivable par rapport à $\theta$ et si le maximum est atteint à l'intérieur de $\Theta$, alors $\hat{\theta}$ vérifie l'équation de vraisemblance :
$$
\frac{\partial l}{\partial \theta}(\hat{\theta}, X) = 0
$$
avec la condition du second ordre :
$$
\frac{\partial^2 l}{\partial \theta^2}(\hat{\theta}, X) < 0
$$
ou, en dimension supérieure, la hessienne doit être définie négative.
\end{proposition}

\begin{theoreme}[Principe d'invariance]
Si $\hat{\theta}$ est un EMV de $\theta$ et $g : \Theta \to \mathbb{R}$ une fonction, alors $g(\hat{\theta})$ est un EMV de $g(\theta)$.
\end{theoreme}

\begin{exemple}
Soit $X_1, \ldots, X_n$ un $n$-échantillon i.i.d. de loi $\mathcal{N}(\theta, 1)$ avec $\theta \in \mathbb{R}$ inconnu.
La vraisemblance est :
$$
L(\theta, X) = \prod_{i=1}^{n} \frac{1}{\sqrt{2\pi}} e^{-\frac{(X_i-\theta)^2}{2}} = (2\pi)^{-n/2} e^{-\frac{1}{2}\sum_{i=1}^{n}(X_i-\theta)^2}
$$
La log-vraisemblance est :
$$
l(\theta, X) = -\frac{n}{2}\ln(2\pi) - \frac{1}{2}\sum_{i=1}^{n}(X_i-\theta)^2
$$
On dérive :
$$
\frac{\partial l}{\partial \theta}(\theta, X) = \sum_{i=1}^{n}(X_i-\theta) = n\bar{X}_n - n\theta
$$
Cette dérivée s'annule en $\theta = \bar{X}_n$. Comme $\frac{\partial^2 l}{\partial \theta^2} = -n < 0$, il s'agit bien d'un maximum.

Donc l'EMV de $\theta$ est : $\hat{\theta} = \bar{X}_n$.

Par le principe d'invariance, l'EMV de $\theta^2$ est $(\bar{X}_n)^2$.
\end{exemple}

\begin{exemple}
Soit $X_1, \ldots, X_n$ un $n$-échantillon i.i.d. de loi exponentielle $\mathcal{E}(\lambda)$ de densité $f_\lambda(x) = \lambda e^{-\lambda x} \mathbb{1}_{x \geq 0}$, avec $\lambda > 0$ inconnu.

La vraisemblance est :
$$
L(\lambda, X) = \prod_{i=1}^{n} \lambda e^{-\lambda X_i} = \lambda^n e^{-\lambda \sum_{i=1}^{n} X_i}
$$
La log-vraisemblance est :
$$
l(\lambda, X) = n\ln(\lambda) - \lambda \sum_{i=1}^{n} X_i
$$
On dérive :
$$
\frac{\partial l}{\partial \lambda}(\lambda, X) = \frac{n}{\lambda} - \sum_{i=1}^{n} X_i
$$
Cette dérivée s'annule en $\lambda = \frac{n}{\sum_{i=1}^{n} X_i} = \frac{1}{\bar{X}_n}$.

Comme $\frac{\partial^2 l}{\partial \lambda^2} = -\frac{n}{\lambda^2} < 0$, il s'agit bien d'un maximum.

Donc l'EMV de $\lambda$ est : $\hat{\lambda} = \frac{1}{\bar{X}_n}$.
\end{exemple}

\begin{remarque}[Avantages et inconvénients]
\textbf{Avantages :}
\begin{itemize}
\item Méthode générale applicable à de nombreux modèles statistiques.
\item Interprétation intuitive : on choisit le paramètre qui rend l'observation la plus probable.
\item Bonnes propriétés asymptotiques (convergence, normalité asymptotique, efficacité).
\item Le principe d'invariance facilite l'estimation de fonctions de paramètres.
\end{itemize}

\textbf{Inconvénients :}
\begin{itemize}
\item L'EMV peut ne pas être unique.
\item L'EMV peut ne pas exister.
\item L'EMV peut être difficile à exhiber (pas de forme explicite).
\item Peut être biaisé en échantillon fini (même si asymptotiquement sans biais).
\end{itemize}
\end{remarque}